% ----------------------------------------------------------------
% Section: Implementation
% ----------------------------------------------------------------
\chapter{IMPLEMENTATION}
\label{ch:implementation}

This chapter describes the implementation of each component in detail, organized by the build pipeline stage in which it is introduced.

\section{U-Boot Patch: PL011-AXI Serial Driver}
\label{sec:impl_uboot_patch}

\subsection{Problem}

The RPi5 firmware Device Tree declares \texttt{serial10} (the default \texttt{stdout-path} target) with \texttt{compatible = "arm,pl011-axi"} \cite{dt_spec_chosen}. U-Boot's PL011 driver (\texttt{drivers/serial/serial\_pl01x.c}) did not include this string in its \texttt{udevice\_id} match table \cite{uboot_serial_pl01x}. As a result, U-Boot couldn't bind the UART node and produced no serial output, making all subsequent debugging impossible.

\subsection{Solution}

Patch \texttt{0001-serial-pl01x-add-arm-pl011-axi-compatible-for-RPi5.patch} adds the \texttt{"arm,pl011-axi"} entry to \texttt{pl01x\_serial\_id[]} with \texttt{TYPE\_PL011}:

\begin{lstlisting}[language=C, caption={Addition to pl01x\_serial\_id[] in serial\_pl01x.c}]
/* RPi5 firmware DT uses "arm,pl011-axi" for serial10 (stdout-path). */
{.compatible = "arm,pl011-axi", .data = TYPE_PL011},
{.compatible = "arm,pl011",     .data = TYPE_PL011},
{.compatible = "arm,pl010",     .data = TYPE_PL010},
\end{lstlisting}

The ARM PL011 TRM (DDI0183G) defines the AXI variant as register-compatible with the standard PL011 \cite{arm_pl011_trm}, making the \texttt{TYPE\_PL011} assignment correct.

\subsection{Risk Assessment}

\textbf{Low risk.} The change only adds a new entry to the match table; existing \texttt{"arm,pl011"} bindings are unaffected. Boards that do not use \texttt{"arm,pl011-axi"} in their Device Tree will not be matched by this entry.

\section{U-Boot AVB Kconfig Configuration}
\label{sec:impl_uboot_avb_kconfig}

Enabling U-Boot AVB requires a specific set of Kconfig symbols. The build stage automates this via \texttt{scripts/config}:

\begin{lstlisting}[language=bash, caption={Kconfig symbols enabled for AVB}]
./scripts/config --enable CONFIG_ANDROID_BOOT_IMAGE
./scripts/config --enable CONFIG_UDP_FUNCTION_FASTBOOT  # selects hidden CONFIG_FASTBOOT
./scripts/config --enable CONFIG_LIBAVB
./scripts/config --enable CONFIG_AVB_VERIFY
./scripts/config --enable CONFIG_CMD_AVB
./scripts/config --set-val CONFIG_FASTBOOT_BUF_ADDR 0x10000000
./scripts/config --set-val CONFIG_AVB_BUF_ADDR 0x10000000
\end{lstlisting}

The buffer address \texttt{0x10000000} (256~MiB) is set explicitly because neither symbol has a default in the upstream \texttt{rpi\_arm64\_defconfig}. The value is chosen to sit above all standard ARM64 boot-image load regions defined in \texttt{board/raspberrypi/rpi/rpi.env}:

\begin{table}[h]
\centering
\caption{RPi5 ARM64 U-Boot load addresses vs.\ AVB buffer placement}
\label{tab:uboot_mem}
\small
\begin{tabular}{llr}
\toprule
\textbf{Region} & \textbf{Symbol / use} & \textbf{Address} \\
\midrule
Kernel Image       & \texttt{kernel\_addr\_r}       & \texttt{0x00080000} \\
Decompression area & \texttt{kernel\_comp\_addr\_r} & \texttt{0x02000000} \\
Device Tree Blob   & \texttt{fdt\_addr\_r}          & \texttt{0x05600000} \\
Ramdisk            & \texttt{ramdisk\_addr\_r}      & \texttt{0x05700000} \\
Ramdisk end (est.) & $\sim$50~MiB max               & \texttt{$\sim$0x08900000} \\
\midrule
\textbf{AVB / Fastboot buffer} & \texttt{CONFIG\_AVB\_BUF\_ADDR} & \textbf{\texttt{0x10000000}} \\
\bottomrule
\end{tabular}
\end{table}

The 119~MiB gap between the estimated ramdisk end and the buffer start ensures no overlap under any realistic ramdisk size. The 112~MiB buffer (\texttt{CONFIG\_AVB\_BUF\_SIZE = 0x7000000}) is the U-Boot default and is sufficient to hold any Android partition image presented to the AVB verification command.

After setting these symbols, \texttt{make olddefconfig} resolves the full dependency graph. The build stage then verifies that all required symbols are present in \texttt{.config} before proceeding; missing symbols raise a descriptive error.

\subsection{Public Key Injection}

The AVB root public key (\texttt{avb\_public\_key.bin}, in AVB packed format) is read and embedded as a C array in \texttt{common/avb\_verify.c}. The build stage uses a regular expression to locate and replace the existing \texttt{avb\_root\_pub[]} array in-place, making the operation idempotent. A warning is printed if signing is enabled but U-Boot isn't being rebuilt, reminding the operator to rebuild U-Boot after any key rotation.

\section{U-Boot Boot Script}
\label{sec:impl_boot_script}

The AVB boot script (\texttt{config/uboot/boot\_avb.cmd}) implements the following sequence:

\begin{enumerate}
    \item Resolve \texttt{avb\_fail\_policy} (default injected at build time via the template placeholder \texttt{\_\_AVB\_FAIL\_POLICY\_\_}).
    \item \texttt{avb init 0} --- attach the AVB device on MMC~0.
    \item \texttt{avb verify} --- verify all descriptors in \texttt{vbmeta.img}. On failure: either reset (fail\_closed) or continue with orange boot state (fail\_open).
    \item \texttt{fatload mmc 0:1} --- load kernel \texttt{Image} and \texttt{ramdisk.img}.
    \item Construct \texttt{bootargs}: append \texttt{avb\_bootargs} (populated by \texttt{avb verify}), SELinux mode, encryption mode, and optional cmdline profile.
    \item \texttt{booti} --- hand off to the Android kernel.
\end{enumerate}

The script is compiled to binary by \texttt{mkimage}:
\begin{lstlisting}[language=bash]
mkimage -C none -A arm64 -T script -d boot_avb.cmd boot_avb.scr
\end{lstlisting}

Template placeholders (\texttt{\_\_SELINUX\_MODE\_\_}, \texttt{\_\_ENCRYPTION\_ARGS\_\_}, \texttt{\_\_CMDLINE\_PROFILE\_ARGS\_\_}) are resolved at build time before \texttt{mkimage} is invoked, keeping the configuration model in the YAML layer rather than hardcoded in the script.

\section{Kernel Patch: Boot-Layer Diagnostic Markers}
\label{sec:impl_kernel_patch}

To measure and verify the progression through the kernel boot sequence, four \texttt{pr\_notice()} markers are injected into \texttt{init/main.c}:

\begin{table}[h]
\centering
\caption{Kernel boot-layer diagnostic markers}
\label{tab:kernel_diag}
\begin{tabular}{cll}
\toprule
\textbf{Layer} & \textbf{Location in init/main.c} & \textbf{Message (prefix: RPi5-DIAG)} \\
\midrule
1 & Entry of \texttt{start\_kernel()} & Layer 1 start\_kernel entry \\
2 & After \texttt{console\_init()} & Layer 2 console\_init complete \\
3 & Before \texttt{do\_initcalls()} & Layer 3 do\_initcalls entry \\
4 & Before \texttt{run\_init\_process()} & Layer 4 kernel\_init userspace exec \\
\bottomrule
\end{tabular}
\end{table}

\texttt{pr\_notice} (priority~5) is visible at the default kernel loglevel and guaranteed to be present in the serial console output. The \texttt{RPi5-DIAG:} prefix makes the markers \texttt{grep}-able in captured logs.

\section{AOSP Patch: AVB-Aware fstab Selection}
\label{sec:impl_avb_fstab}

The upstream RPi5 device tree does not provide an AVB-aware \texttt{fstab}.

Patch \texttt{0001-rpi5-add-avb-aware-fstab-selection.patch} adds:

\begin{itemize}
    \item \texttt{RPI5\_ENABLE\_AVB} build toggle in \texttt{BoardConfig.mk}.
    \item \texttt{androidboot.vbmeta.device=/dev/block/by-name/vbmeta} kernel command line argument, so first-stage init can locate the VBMeta partition by GPT label.
    \item \texttt{ramdisk/fstab.rpi5.avb}: AVB-aware fstab with \texttt{avb=vbmeta} on both \texttt{/system} and \texttt{/vendor}.
    \item Conditional selection in \texttt{device.mk}: when \texttt{RPI5\_ENABLE\_AVB=true}, the AVB fstab is used in both ramdisk and vendor partitions.
\end{itemize}

\begin{lstlisting}[caption={fstab.rpi5.avb --- system and vendor entries}]
/dev/block/by-name/system  /system  ext4  ro,barrier=1
    wait,avb=vbmeta,first_stage_mount
/dev/block/by-name/vendor  /vendor  ext4  ro,barrier=1
    wait,avb=vbmeta,first_stage_mount
\end{lstlisting}

The \texttt{avb=vbmeta} flag instructs first-stage init to look up the VBMeta descriptor for each partition before mounting it.

\section{AOSP Patch: SELinux Enforcing Mode}
\label{sec:impl_selinux}

The upstream RPi5 AOSP device tree ships SELinux in Permissive mode, where policy violations are logged but not enforced. The objective was to enable SELinux Enforcing mode and confirm that the platform's base policy is compatible with RPi5 bring-up.

\subsection{Policy Structure}

\begin{sloppypar}
Android precompiles the SELinux policy at build time into \path{/vendor/etc/selinux/precompiled_sepolicy}. The policy is a binary merge of the platform policy (\texttt{PLAT\_SEPOLICY}), vendor policy (\texttt{BOARD\_SEPOLICY}), and their mapping files. At boot, \texttt{init} validates the precompiled policy against SHA-256 hashes stored in \path{plat_sepolicy_and_mapping.sha256} before loading it into the kernel via \path{security_load_policy()} \cite{selinux_build_precompiled}.
\end{sloppypar}

The RPi5 device tree does not define device-specific policy modules. It inherits the base AOSP platform policy unchanged, which is sufficient for the subsystems under evaluation (AVB, FBE, boot diagnostics). No custom \texttt{.te} policy files or file-context entries were added beyond those introduced by the \texttt{e2fsdroid} fix described in Section~\ref{sec:impl_fbe_selinux}.

\subsection{Android 16 Enforcing Behaviour}

In Android~16 (\texttt{android-16.0.0\_r4}), SELinux transitions to Enforcing mode unconditionally at policy load time. The \texttt{androidboot.selinux=permissive} kernel argument, honoured in earlier Android releases, is ignored under this build configuration. \texttt{init}'s \texttt{IsEnforcing()} calls \texttt{security\_getenforce()} after loading the precompiled policy and does not re-read the kernel command line permissive flag. There is therefore no safe runtime fallback to Permissive mode: any policy violation preventing a required service from starting will cause a boot failure.

\subsection{Configuration}

SELinux mode is controlled by the \texttt{boot.selinux\_mode} key in the rpi5 build config YAML (e.g.\ \texttt{config/rpi5\_uboot\_aosp\_signed.yaml}):

\begin{lstlisting}[language=bash, caption={boot.selinux\_mode in rpi5\_uboot\_aosp\_signed.yaml}]
boot:
  # Options: permissive | enforcing
  selinux_mode: permissive   # set to enforcing for production
\end{lstlisting}

During the build stage, \texttt{stages/build.py} renders the \texttt{\_\_SELINUX\_MODE\_\_} placeholder in the boot script template (\texttt{config/uboot/boot\_avb.cmd}) with this value:

\begin{lstlisting}[language=bash, caption={androidboot.selinux injected into bootargs in boot\_avb.cmd}]
setenv bootargs "${bootargs} ... androidboot.selinux=__SELINUX_MODE__ ..."
\end{lstlisting}

The mode can also be overridden per-run via the \texttt{--selinux-mode} CLI flag without modifying the YAML. To enable SELinux Enforcing mode, set \texttt{boot.selinux\_mode: enforcing} in the config YAML before running the build stage.

\subsection{AVC Denial Handling}

During bring-up, SELinux enforcement failures appear as AVC (Access Vector Cache) denial messages in both \texttt{dmesg} and \texttt{logcat}:

\begin{lstlisting}
avc: denied { <permission> } for comm="<process>"
    scontext=<source_context> tcontext=<target_context>
    tclass=<object_class> permissive=0
\end{lstlisting}

The \texttt{permissive=0} field confirms the denial was enforced rather than just logged. The \texttt{rpi5\_boot\_diag.sh} service captures the SELinux enforcement mode (\texttt{ro.boot.selinux}, \texttt{getenforce}) as part of its post-boot state report, providing per-boot confirmation of the enforcement posture.

\subsection{Serial Log Evidence}
\label{sec:impl_selinux_logs}

\subsubsection*{Non-Secure Device (\texttt{selinux\_mode: permissive})}

The signed build with \texttt{selinux\_mode: permissive} produced the following kernel command line and property state:

\begin{lstlisting}[caption={Kernel command line --- permissive build}]
Kernel command line: stack_depot_disable=on cgroup_disable=pressure
  root=/dev/ram0 rootwait androidboot.hardware=rpi5
  androidboot.selinux=permissive ignore_loglevel loglevel=7
  initcall_debug androidboot.verifiedbootstate=green
  androidboot.vbmeta.device_state=locked
  androidboot.serialno=6113ca6606956aab
\end{lstlisting}

\begin{lstlisting}[caption={Selected getprop output --- permissive build}]
[ro.boot.selinux]:           [permissive]
[ro.boot.verifiedbootstate]: [green]
[ro.crypto.state]:           [unsupported]
\end{lstlisting}

\subsubsection*{Secure Device (\texttt{selinux\_mode: enforcing})}

The signed build with \texttt{selinux\_mode: enforcing} produced:

\begin{lstlisting}[caption={Kernel command line --- enforcing build}]
Kernel command line: stack_depot_disable=on cgroup_disable=pressure
  root=/dev/ram0 rootwait androidboot.hardware=rpi5
  androidboot.selinux=enforcing ignore_loglevel loglevel=7
  initcall_debug androidboot.verifiedbootstate=orange
  androidboot.vbmeta.device_state=unlocked
  androidboot.serialno=6113ca6606956aab
\end{lstlisting}

\begin{lstlisting}[caption={Selected getprop output --- enforcing build}]
[ro.boot.selinux]:           [enforcing]
[ro.boot.verifiedbootstate]: [orange]
[ro.crypto.state]:           [unsupported]
\end{lstlisting}

\begin{lstlisting}[caption={Policy transition and first AVC denials --- enforcing build}]
[23.245] audit: enforcing=1 old_enforcing=0 lsm=selinux
[23.407] avc: denied { set } for
           property=ro.boot.audio.tinyalsa.simulate_input
           scontext=u:r:vendor_init:s0
           tcontext=u:object_r:bootloader_prop:s0
           tclass=property_service permissive=0
[23.451] avc: denied { set } for
           property=persist.bluetooth.a2dp_aac.vbr_supported
           scontext=u:r:vendor_init:s0
           tcontext=u:object_r:bluetooth_prop:s0
           tclass=property_service permissive=0
[23.495] avc: denied { set } for
           property=ro.lockscreen.disable.default
           scontext=u:r:vendor_init:s0
           tcontext=u:object_r:default_prop:s0
           tclass=property_service permissive=0
[23.683] avc: denied { search } for comm="ueventd" name="/"
           dev="mmcblk0p7" scontext=u:r:ueventd:s0
           tcontext=u:object_r:unlabeled:s0
           tclass=dir permissive=0
\end{lstlisting}

The \texttt{permissive=0} field in every entry confirms enforcement was active. The denials are RPi5-specific policy gaps (vendor properties, unlabelled \texttt{/metadata} root inode) that require device-specific policy extensions for a production deployment.

\section{AOSP Patch: File-Based Encryption}
\label{sec:impl_fbe}

Patch \texttt{0002-rpi5-add-fbe-file-based-encryption.patch} enables FBE on \texttt{userdata} \cite{android_fbe}:

\begin{itemize}
    \item \texttt{RPI5\_ENABLE\_FBE} toggle extends the three-way fstab selection: FBE $>$ AVB-only $>$ base.
    \item \texttt{ramdisk/fstab.rpi5.fbe}: extends the AVB fstab with FBE flags on \texttt{/data} and a \texttt{/metadata} entry.
    \item \texttt{ro.crypto.type=file} product property signals the crypto stack.
    \item \texttt{androidboot.encryption\_mode=fbe} is appended to the kernel command line.
\end{itemize}

The critical \texttt{fstab} entry for \texttt{userdata}:

\begin{lstlisting}[caption={fstab.rpi5.fbe --- userdata entry (FBE + metadata encryption)}]
/dev/block/by-name/userdata  /data  ext4
    nodev,noatime,nosuid,errors=panic
    wait,check,formattable,quota,
    fileencryption=aes-256-xts:aes-256-cts:v2,
    keydirectory=/metadata/vold/metadata_encryption
\end{lstlisting}

The \texttt{fileencryption=aes-256-xts:aes-256-cts:v2} flag selects AES-256-XTS for file content and AES-256-CTS for filenames at policy version~2. The \texttt{keydirectory=} flag activates metadata encryption: \texttt{vold} stores the wrapped key in the \texttt{/metadata} partition, which is itself not encrypted, ensuring key material survives a factory reset without persisting in \texttt{/data} \cite{android_data_encryption}.

\section{FBE Bring-up: Metadata Partition SELinux Boot Loop}
\label{sec:impl_fbe_selinux}

\subsection{Symptom}

After flashing a signed SD card image with FBE enabled, the device entered a repeating reboot loop. Each cycle completed in approximately 9--10 seconds, with AVB verification passing on every cycle:

\begin{lstlisting}[caption={Reboot loop symptom on UART console}]
AVB: Verification PASSED
...
[2.8s] vold: Failed to create /metadata/vold: Permission denied
[2.8s] vold: read_key failed in mountFstab
[9.4s] reboot: Restarting system with command 'bootloader'
\end{lstlisting}

Because AVB verification passed, the signing pipeline was confirmed correct. The failure was inside the FBE initialization path.

\subsection{Diagnosis: Two Concurrent Root Causes}

Full debug visibility was obtained by rebuilding \texttt{boot\_avb.scr} with \texttt{--cmdline-profile debug} (\texttt{ignore\_loglevel loglevel=7 initcall\_debug}), which exposed the complete init sequence on the serial console.

\subsubsection{Root Cause A --- SELinux Enforcing Despite Permissive Kernel Argument}

The bootargs carried \texttt{androidboot.selinux=permissive}, but the serial log showed:

\begin{lstlisting}
audit: enforcing=1 old_enforcing=0
\end{lstlisting}

\begin{sloppypar}
All subsequent AVC denials showed \texttt{permissive=0} (enforced). In Android~16 (\path{android-16.0.0_r4}), init's \texttt{IsEnforcing()} function ignores the command-line permissive flag under this build configuration --- SELinux transitions to enforcing unconditionally at policy load time.
\end{sloppypar}

\subsubsection{Root Cause B --- Metadata Partition Root Inode is \texttt{unlabeled}}

The SD card build stage created the \texttt{metadata} filesystem with bare \texttt{mkfs.ext4}, which produces a filesystem with no SELinux extended attributes. The partition root directory inode therefore carried the default context \texttt{u:object\_r:unlabeled:s0}.

The compiled \texttt{plat\_file\_contexts} already contained the correct mapping:

\begin{lstlisting}
/metadata(/.*)?       u:object_r:metadata_file:s0
/metadata/vold(/.*)?  u:object_r:vold_metadata_file:s0
\end{lstlisting}

The label was simply never applied to the inode. When vold (context \texttt{u:r:vold:s0}) attempted to create \texttt{/metadata/vold/}, SELinux denied it:

\begin{lstlisting}
avc: denied { write } for comm="binder:344_2" name="/"
    scontext=u:r:vold:s0 tcontext=u:object_r:unlabeled:s0
    tclass=dir permissive=0
\end{lstlisting}

Without \texttt{/metadata/vold/}, \texttt{vold} couldn't store the FBE master key, which is required by the \texttt{keydirectory=} fstab flag before \texttt{/data} can be mounted.

\subsection{Fix: Replace \texttt{mkfs.ext4} with \texttt{e2fsdroid}}

The fix replaced the bare \texttt{mkfs.ext4} call in the SD card stage with \texttt{e2fsdroid}, the AOSP-provided tool that creates ext4 filesystems with SELinux extended attributes pre-populated from the compiled \texttt{file\_contexts} map:

\begin{lstlisting}[language=bash, caption={e2fsdroid invocation for labelled metadata partition}]
e2fsdroid -e -T 0 \
    -S ${SELINUX_FC}  \   # compiled file_contexts
    -L /metadata       \   # mount point label root
    ${METADATA_IMAGE}
\end{lstlisting}

With \texttt{e2fsdroid}, the metadata partition root inode is created with \path{u:object_r:metadata_file:s0}, vold can create \texttt{/metadata/vold/}, and FBE initialisation succeeds.

\subsection{Security Significance}

This incident demonstrates the interdependency between the MAC layer (SELinux) and the encryption layer (FBE): a correct FBE \texttt{fstab} configuration fails silently if the partition on which the key material must be stored lacks the correct SELinux label. The root cause was not a security vulnerability but a toolchain gap --- \texttt{mkfs.ext4} does not apply SELinux labels, while \texttt{e2fsdroid} does. For a verified boot stack, where every layer must be correctly labelled for the MAC policy to permit operations, using the AOSP-provided image creation tools is mandatory.

\section{AOSP Patch: Boot Diagnostics Service}
\label{sec:impl_diagnostics}
Patch \texttt{0003-rpi5-add-boot-diagnostics.patch} adds a post-boot diagnostic service:

\begin{sloppypar}
\begin{description}
    \item[\texttt{init.rpi5.diag.rc}] Defines a \texttt{oneshot} service triggered by \path{on property:sys.boot_completed=1}. The \texttt{console} keyword routes \texttt{stdout} to \texttt{/dev/console}, making output appear on the UART without ADB.
    \item[\texttt{rpi5\_boot\_diag.sh}] Collects and prints: build info, AVB verified boot state (\path{ro.boot.verifiedbootstate}), VBMeta digest, crypto state (\texttt{ro.crypto.state}, \texttt{ro.crypto.type}), FBE properties, \texttt{vold} status, \texttt{fstab}, mount table, \texttt{dm-*} block devices, and SELinux enforce mode. Output is simultaneously written to \path{/data/local/tmp/rpi5_boot_diag.log}.
\end{description}
\end{sloppypar}

This service provides an unambiguous, per-boot record of the security state of the device without requiring ADB connectivity, which is particularly useful during bring-up when network and USB may not yet be functional.

\section{AVB Signing Stage}
\label{sec:impl_sign}

The sign stage (\texttt{src/meta\_rpi5\_secure\_aosp/stages/sign.py}) performs:

\begin{enumerate}
    \item For each of \texttt{system.img} and \texttt{vendor.img}:
    \begin{enumerate}
        \item Copy to \texttt{.signed.img}.
        \item Calculate the maximum payload size for an AVB hashtree footer via \texttt{avbtool calculate\_max\_image\_size}.
        \item If the image exceeds this size: \texttt{e2fsck -fy} $\to$ \texttt{resize2fs} $\to$ \texttt{truncate} to the target size.
        \item Append the AVB hashtree footer via \texttt{avbtool add\_hashtree\_footer} (algorithm: SHA256/RSA-4096, key: \texttt{avb\_private\_key.pem}).
    \end{enumerate}
    \item \begin{sloppypar}Generate \texttt{vbmeta.img} via \path{avbtool make_vbmeta_image}, referencing the signed \path{system.signed.img} and \path{vendor.signed.img}
    \cite{avb2_readme, avbtool_github}.\end{sloppypar}
\end{enumerate}

The \texttt{AvbTool} wrapper class in \texttt{utils/avb.py} abstracts the \texttt{avbtool} invocations and resolves the tool path from the AOSP source tree.

\section{Build Orchestrator CLI}
\label{sec:impl_cli}

The \texttt{rpi5-build} CLI (\texttt{src/meta\_rpi5\_secure\_aosp/main.py}) provides:

\begin{itemize}
    \item Stage selection: \texttt{--stage \{all|sync|patch|build|sdcard\}}
    \item Component selection: \texttt{--code \{all|kernel|uboot|aosp\}}
    \item Per-run security overrides: \texttt{--build-variant}, \texttt{--selinux-mode}, \texttt{--boot-state-override}, \texttt{--avb-fail-policy}, \texttt{--encryption-mode}
    \item An explicit \texttt{--allow-insecure-boot-state} guard that must be provided before any insecure option (orange state, fail\_open) is accepted.
\end{itemize}

The \texttt{BuildContext} dataclass is the single shared state object threaded through all stage modules, eliminating global state and making each stage independently testable.
